\chapter{Testing and Deploying Low-Latency Systems}\label{ch:ll:testing-and-deploying-low-latency-systems}

A release of a trading system can pass every functional test and still be worse than the one it replaces: a log statement
that moved inside the hot loop, a container that now allocates, a lock taken on a path that used to be free. None of these
changes an output, so no functional test fails; each adds tens or hundreds of nanoseconds at the 99th percentile, which is
where a latency-sensitive strategy lives. Catching them takes a different kind of test, one that measures, on a machine that
is never perfectly quiet, and decides whether a difference is real. This chapter assembles the tests a low-latency system
needs, from the ones that check what it computes (properties, differences, random inputs, the venue's own scripts) to the one
that checks how fast it computes, and ends with the way a release reaches production: gradually, reversibly, and only onto
machines that have been checked.

\section{The test pyramid for a trading system}

\begin{definition}[Property-based test, differential test, fuzz testing]\label{def:ll:testing-and-deploying-low-latency-systems:tests}
A \emph{property-based test}\index{property-based test} generates many random inputs and checks that a stated property holds
for every one of them, rather than comparing one input with one expected output. A \emph{differential
test}\index{differential test} runs two implementations of the same specification on the same inputs and checks that their
outputs are identical. \emph{Fuzz testing}\index{fuzz testing} feeds a program large numbers of malformed or mutated inputs to
find those that make it crash, hang or violate its invariants.
\end{definition}

The components of this book have been tested in these ways all along. The gateway of chapter~21 is checked by properties over
random interleavings against the matching engine (no order fills more than it asked for, no position exceeds the limit); the
\omterm{def:ll:the-order-book-builder:builder}{book builder} of chapter~19 and the engines of chapters~20 to 24 by \omterm{def:ll:testing-and-deploying-low-latency-systems:tests}{differential tests}, since the Python reference, the C++ and
the Rust implementations must agree on every fixture, byte for byte or hash for hash. Property-based testing is the idea of
QuickCheck (Claessen and Hughes, 2000): the property is the specification, and the generator explores inputs no one would have
written by hand. \Cref{fig:ll:testing-and-deploying-low-latency-systems:pyramid} puts the kinds of test in the order in which a
release meets them: many cheap ones at the bottom, run on every change, and a few expensive ones at the top, run before
production.

\begin{omfigure}
\begin{tikzpicture}[font=\scriptsize]
\foreach \y/\w/\t/\c in {0/11/{unit, property and differential tests: every change, seconds}/omDef,
  0.9/9.2/{fuzzing with sanitisers: every change, minutes}/omDef,
  1.8/7.4/{simulator-backed scenarios and replays of recorded days}/omMeth,
  2.7/5.6/{certification scripts against the venue's test system}/omMeth,
  3.6/3.8/{performance gate: interleaved runs}/omThm,
  4.5/2.2/{staged rollout}/omThm} {
  \fill[\c!18,draw=\c!60] (-\w/2,\y) rectangle (\w/2,\y+0.8);
  \node[align=center] at (0,\y+0.4) {\t};
}
\draw[->] (6.3,0.2) -- node[right,align=left]{fewer, slower,\\ closer to production} (6.3,5.1);
\end{tikzpicture}

\omcaption{The tests a low-latency trading system meets on its way to production. The lower layers check what it computes and
run on every change; the upper ones check how it behaves against the venue and how fast it is, and run on candidates for
release.}\label{fig:ll:testing-and-deploying-low-latency-systems:pyramid}
\end{omfigure}

\section{Tests backed by the exchange simulator}

Unit tests stop at the process boundary; the bugs that cost money live at it: a report that arrives before the acknowledgement
of its order, a cancel that crosses a fill, a session that drops in the middle of a burst, a venue that rejects a message the
firm thought valid. Book 10's simulator makes these reproducible: the same seed gives the same interleaving, a scripted fault
(a dropped session, a halt, a pause) happens at the same moment, and the matching engine applies the venue's rules exactly.
Chapter~21's property tests and chapter~24's \omterm{def:ll:resilience:failover}{failover} table are simulator-backed tests in this sense: a scenario, a set of
seeds or cut points, and properties checked after each run. Replays of recorded days (chapter~23) complete them: they check
that a new build takes the decisions the old one took, except where a change was intended.

Random inputs reach places scenarios do not. \Cref{lst:ll:testing-and-deploying-low-latency-systems:trusting} is a checksum
reader with a classic bug: it believes the message's BodyLength field and reads the checksum where the field says, without
checking that the message is that long. The chapter's fuzzer mutates the golden FIX messages of chapter~15 (a byte flipped,
a digit changed, a field duplicated, a tail cut off) and copies each into a buffer of exactly its size, so that
AddressSanitizer sees any read past it. Against this reader it stops at the first truncated message: within 12 inputs for
each of ten seeds. Against the firm's own parser it ran a million inputs without a bad read. It did find a real bug there,
through an invariant rather than a crash: two mutations, a tag 11 turned into 10 and a sequence number raised by one, left the
checksum unchanged, and the parser accepted a message with a CheckSum field in the middle of its body. The Python, C++ and Rust
parsers of \texttt{firm.fixengine} now refuse it, and a test keeps the message.

\omcode{code/low-latency/25-testing-and-deploying-low-latency-systems/cpp/ll_fuzz.hpp}{67}{78}{A checksum reader that trusts
BodyLength: the fuzzer's first truncated message makes it read past the buffer.}\label{lst:ll:testing-and-deploying-low-latency-systems:trusting}

\section{Exchange certification and conformance}

\begin{definition}[Exchange certification]\label{def:ll:testing-and-deploying-low-latency-systems:cert}
\emph{Exchange certification}\index{exchange certification} is the process by which a venue checks, before granting access to
production, that a client's system handles its protocols correctly: the client runs a script of scenarios against the venue's
test system, which verifies each message and response.
\end{definition}

Venues certify the systems that connect to them, and regulators ask firms to test their conformance with each venue before
trading and after every material change (the dated box). A certification script is a list of steps, each a message and the
responses it must produce: an order accepted, a replace, a cancel, a cancel that comes too late, a reused identifier, a price
off the tick grid, an immediate-or-cancel order with nothing to trade, a post-only order that would cross, a \omterm{def:ll:pre-trade-risk-and-kill-switches:mass}{mass cancel}, a
disconnect that cancels resting orders. \texttt{firm.perfgate}'s runner plays such a script against Book 10's matching engine
(\cref{lst:ll:testing-and-deploying-low-latency-systems:certify}); the build's seventeen steps pass, and a script that expects
the wrong behaviour of a single step fails at that step and no other. The same runner, pointed at a new version of the
simulator or of the firm's gateway, is a conformance test in the regulator's sense.

\omcode{code/firm/perfgate/firm_perfgate.py}{121}{142}{The certification runner: each step sends one message and compares the
responses, in order, with the script's.}\label{lst:ll:testing-and-deploying-low-latency-systems:certify}

\begin{dated}{2026-09}{Conformance testing and certification}\label{dat:ll:testing-and-deploying-low-latency-systems:cert}
In the European Union (consulted September 2026), Commission Delegated Regulation (EU) 2017/589 requires firms to establish
methodologies to develop and test algorithmic trading systems before deployment or substantial update; to test their
conformance with a trading venue's system when becoming a member, when connecting through sponsored access for the first time,
when the venue's systems change materially and before deploying or materially updating an algorithm, verifying that it
interacts with the venue's matching logic as intended; and to test in an environment separated from production. CME Group
requires client systems that route orders to or process market data from its Globex platform to be certified with AutoCert+,
its automated testing tool.
\end{dated}

\section{Performance regression gates}

\begin{definition}[Performance regression gate]\label{def:ll:testing-and-deploying-low-latency-systems:gate}
A \emph{performance regression gate}\index{performance regression gate} is a test that runs a benchmark on the current build
and on a candidate, on the same machine, and refuses the candidate if a latency statistic is worse by more than a stated
tolerance with a stated confidence.
\end{definition}

A single benchmark run says little, even on an otherwise idle laptop. \Cref{fig:ll:testing-and-deploying-low-latency-systems:runs} shows 800
runs of the gate's benchmark (the strategy engine on the small-tick day of chapter~19, timed per event), alternating between the
current build and one with a planted regression. The 99th percentile of a run moves by tens of percent from one run to the
next with the operating system's and the hypervisor's own activity; the regression, which adds about 18\% at the median of
the runs, is invisible in
any single pair. The gate is built for this noise. It interleaves the runs (A, B, B, A, \dots), so that a slow drift of the
machine affects both builds alike; it takes one statistic per run, the 99th percentile, and compares their medians across runs;
it decides with a permutation test, which assumes nothing about the distribution, and reports a bootstrap interval of the rise
(\cref{lst:ll:testing-and-deploying-low-latency-systems:compare}); and it flags a regression only if the rise is both significant
and larger than a tolerance of 2\%, so that a real but negligible change does not block a release.

\omcode{code/firm/perfgate/firm_perfgate.py}{62}{83}{The gate's decision: the rise of the median 99th percentile, a one-sided
permutation test, a bootstrap interval, and a tolerance.}\label{lst:ll:testing-and-deploying-low-latency-systems:compare}

\begin{omfigure}
\begin{tikzpicture}
\begin{axis}[width=11cm,height=5.4cm,xmin=0,xmax=800,ymin=150,ymax=700,xlabel={run (interleaved: A B B A \dots)},
  ylabel={99th percentile of the run (ns)},tick label style={font=\small},label style={font=\small},grid=major,
  grid style={gray!25},legend pos=north east,legend style={font=\scriptsize},table/col sep=comma]
\addplot[only marks,mark=*,mark size=0.7pt,omDef,restrict expr to domain={\thisrow{isb}}{0:0}]
  table[x=run,y=p99_ns]{figdata/low-latency/25-testing-and-deploying-low-latency-systems/gate_runs_plot.csv};
\addlegendentry{A: current build}
\addplot[only marks,mark=*,mark size=0.7pt,omThm,restrict expr to domain={\thisrow{isb}}{1:1}]
  table[x=run,y=p99_ns]{figdata/low-latency/25-testing-and-deploying-low-latency-systems/gate_runs_plot.csv};
\addlegendentry{B: planted regression}
\end{axis}
\end{tikzpicture}

\omcaption{The 99th percentile of each of 800 runs of the gate's benchmark (18\,000 timed events each), alternating between
the current build and one with a planted regression, on a laptop (Intel Core Ultra 7 155H, WSL2), the machine otherwise idle.
Data: \texttt{bench\_gate.py}.}\label{fig:ll:testing-and-deploying-low-latency-systems:runs}
\end{omfigure}

How many runs does the gate need? \Cref{fig:ll:testing-and-deploying-low-latency-systems:power} answers from the machine's own
noise. It draws gates of $n$ runs per side from the measured pools: two disjoint draws from the current build's runs give the
false-alarm rate, which stays near the test's level of 5\% (3 to 9\% over a hundred gates); two draws of which one is shifted
by 5\% of the median give the detection rate of a 5\% rise; and the planted regression against the current build gives its
detection rate. The planted regression, at about 18\%, is caught 99\% of the time with 40 runs per side. A 5\% rise needs far
more: 73\% at 120 runs, 84\% at 160 and 91\% at 200, so about 200 runs per side for 90\%; the normal approximation of
\texttt{runs\_needed}, from the runs' spread of about 14\%, says about 220. At some 30 milliseconds a run, that is about twelve
seconds of benchmark per gate on this laptop, and on
a production-like machine with isolated cores (chapter~13), whose runs vary less, it is far fewer.

\begin{omfigure}
\begin{tikzpicture}
\begin{axis}[width=8.5cm,height=5.6cm,xmode=log,xmin=4,xmax=250,ymin=0,ymax=1.05,
  xlabel={runs per side},ylabel={fraction of gates that flag B},tick label style={font=\small},label style={font=\small},
  grid=major,grid style={gray!25},legend pos=outer north east,legend style={font=\scriptsize,cells={anchor=west}},
  table/col sep=comma,xtick={5,10,20,40,80,200},xticklabels={5,10,20,40,80,200}]
\addplot[omThm,thick,mark=*] table[x=runs,y=detect_planted]{figdata/low-latency/25-testing-and-deploying-low-latency-systems/gate.csv};
\addlegendentry{planted regression (about 18\%)}
\addplot[omMeth,thick,mark=square*] table[x=runs,y=detect_5pct]{figdata/low-latency/25-testing-and-deploying-low-latency-systems/gate.csv};
\addlegendentry{a 5\% rise}
\addplot[omDef,thick,mark=triangle*] table[x=runs,y=false_alarm]{figdata/low-latency/25-testing-and-deploying-low-latency-systems/gate.csv};
\addlegendentry{no change (false alarms)}
\addplot[gray,dashed] coordinates {(4,0.9) (250,0.9)};
\end{axis}
\end{tikzpicture}

\omcaption{The gate's detection and false-alarm rates against the number of runs per side, from 100 gates drawn from the
measured runs of \cref{fig:ll:testing-and-deploying-low-latency-systems:runs} for each point; the dashed line is 90\% power.
Data: \texttt{fig\_gate.py} (deterministic, from the measured pools).}\label{fig:ll:testing-and-deploying-low-latency-systems:power}
\end{omfigure}

\section{Staged rollouts and rollback}

A release that has passed the gate still meets production gradually: the canary deployment and staged rollout of One Quant
Book~7 (chapter~21), with paper trading where the strategy allows it. Low-latency systems add three habits. The host is checked
before the process starts: \texttt{firm.perfgate.preflight} runs the tuning audit of chapter~13 against the core placement plan
and stops the deployment on any failure, because a strategy tested on isolated cores and deployed on a host whose interrupts land
on them is not the strategy that was tested (this laptop fails 7 of the audit's 11 rules, which is why its measurements carry
the words ``no isolated cores''). The new build runs first next to the old one, on the same inputs, in shadow: the replay of
chapter~23 made live, with its outputs compared and discarded. And rollback is a deployment like any other, prepared and
tested before it is needed, since the moment it is needed is the moment the team has the least time; the \omterm{def:ll:resilience:failover}{failover} of
chapter~24 is the same lesson for machines.

\section{Tutorial: fuzz, certify, gate}\label{tut:ll:testing-and-deploying-low-latency-systems:tut}

\begin{tutorial}
\textbf{Goal.} Fuzz the FIX parser under the sanitisers, run a certification script against the simulator, and gate a planted
regression on this laptop. \textbf{End state:} the fuzzing results, a passing certification, and
\cref{fig:ll:testing-and-deploying-low-latency-systems:power}.
\begin{tutsteps}
\item \textbf{Fuzz.} \texttt{ll\_fuzz\_test.cpp} runs 200\,000 mutations through the firm's parser and checks that every
accepted message lies within its buffer and carries a correct checksum; \texttt{python bench\_fuzz.py} builds the fuzzer with
AddressSanitizer and UndefinedBehaviorSanitizer and runs it against both parsers.
\item \textbf{Certify.} \texttt{firm\_perfgate.certify(load\_script(\textquotedbl{}data/cert\_script.json\textquotedbl{}))}
plays the seventeen steps against Book 10's engine.
\item \textbf{Gate.} \texttt{python bench\_gate.py} measures 400 interleaved pairs; \texttt{python fig\_gate.py} turns them into
detection and false-alarm rates.
\item \textbf{Preflight.} \texttt{firm\_perfgate.preflight} on the tuning audit's two fixture hosts: the tuned one passes, the
mistuned one does not.
\end{tutsteps}
\textbf{What to change next.} Gate on the median of the runs' medians instead of their 99th percentiles and compare the number
of runs needed; add a step to the certification script for the venue's throttle.
\end{tutorial}

\section{Build: the performance gate}\label{bld:ll:testing-and-deploying-low-latency-systems:perfgate}

\begin{build}
\bfield{Purpose} The last checks before production for every component of the book: its speed against the current build, its
behaviour against the venue's script, and the host it will run on.
\bfield{Interface} Python \texttt{firm\_perfgate}: \texttt{interleave}, \texttt{run\_pairs(cmd\_a, cmd\_b, n\_pairs, parse)},
\texttt{compare(a, b, tolerance, alpha)} returning a \texttt{Verdict} with \texttt{rise}, \texttt{p\_value}, \texttt{low},
\texttt{high}, \texttt{regression} and \texttt{report()}; \texttt{power}, \texttt{runs\_needed}; \texttt{certify(script)},
\texttt{load\_script}; \texttt{preflight(root, plan)}.
\bfield{Rules} Runs interleaved, never all of A then all of B; one statistic per run; a permutation test and a tolerance; the
number of runs set from the machine's measured noise; no deployment on a host that fails the audit.
\bfield{Acceptance tests} \texttt{code/firm/perfgate/}: a clear rise flagged and noise or a rise under the tolerance not; false
alarms near the test's level on synthetic noise; power that grows with the number of runs; the certification script passing,
and failing at exactly the altered step; the preflight passing the tuned fixture host and failing the mistuned one.
\bfield{Stretch} A gate on the whole latency distribution (a two-sample test on pooled samples, with care for dependence within
runs); certification steps for sessions and throttles; the gate in continuous integration on a dedicated machine.
\end{build}

\begin{omsources}
\item K.~Claessen and J.~Hughes, ``QuickCheck: a lightweight tool for random testing of Haskell programs'', \emph{ICFP}, 2000.
\item Commission Delegated Regulation (EU) 2017/589 (RTS 6), articles 5 to 7.
\item CME Group Client Systems Wiki, \emph{Client Application Testing and Certification}.
\end{omsources}

\section{Exercises}

\begin{exercise}[$\star$]\label{exo:ll:testing-and-deploying-low-latency-systems:1}
State two properties that a test over random interleavings should check for the \omterm{def:ll:order-gateway-and-order-management:gateway}{order gateway} of chapter~21, and one that it
cannot check without the venue's \omterm{def:ll:order-gateway-and-order-management:dropcopy}{drop copy}.
\end{exercise}

\begin{exercise}[$\star$]\label{exo:ll:testing-and-deploying-low-latency-systems:2}
Why does the fuzzer copy each input into a buffer of exactly its size before parsing it?
\end{exercise}

\begin{exercise}[$\star$]\label{exo:ll:testing-and-deploying-low-latency-systems:3}
The runs of A and B are interleaved A, B, B, A. What goes wrong if all the runs of A come first, then all the runs of B?
\end{exercise}

\begin{exercise}[$\star\star$]\label{exo:ll:testing-and-deploying-low-latency-systems:4}
With runs whose 99th percentiles vary by 14\%, how many runs per side does the normal approximation ask for to detect a 10\%
rise with 90\% power, and a 2.5\% rise?
\end{exercise}

\begin{exercise}[$\star\star$]\label{exo:ll:testing-and-deploying-low-latency-systems:5}
A gate runs 20 times a day at a level of 5\%. How many false alarms should the team expect in a month of 21 trading days, and
what does the tolerance of 2\% change?
\end{exercise}

\begin{exercise}[$\star\star$]\label{exo:ll:testing-and-deploying-low-latency-systems:6}
Two mutations changed a tag from 11 to 10 and a sequence number from 2 to 3, and the checksum did not change. Explain why, and
say what the FIX checksum does and does not protect against.
\end{exercise}

\begin{exercise}[$\star\star\star$]\label{exo:ll:testing-and-deploying-low-latency-systems:7}
\emph{Coding.} Add certification steps for the venue's throttle: configure the engine with a rate, send a burst, and check the
rejections with reason \texttt{T}. What does the step need that the others did not?
\end{exercise}

\begin{exercise}[$\star\star\star$]\label{exo:ll:testing-and-deploying-low-latency-systems:8}
\emph{Find the flaw.} ``Our performance test runs the benchmark once on the release candidate and fails if the 99th percentile
is more than 5\% above last week's number, which we keep in a file.''
\end{exercise}

\section{Problem: A Few Hundred Nanoseconds}

\begin{problem}[{Weekend problem --- how many runs a gate needs on a noisy machine}]\label{pb:ll:testing-and-deploying-low-latency-systems:1}
Use \cref{fig:ll:testing-and-deploying-low-latency-systems:runs,fig:ll:testing-and-deploying-low-latency-systems:power}
(\texttt{measured\_gate\_runs.csv}, \texttt{gate.csv}, \texttt{gate\_summary.csv}).

\medskip\noindent\textbf{Part I --- The noise.}
\begin{enumerate}
\item What is the median 99th percentile of the current build's runs, and how much do runs vary (the robust coefficient of
variation)?
\item What does the planted regression add, at the median of the runs?
\item Why is the 99th percentile of a run so much noisier than its median?
\item Why is a single A/B pair useless on this machine?
\end{enumerate}

\medskip\noindent\textbf{Part II --- The gate.}
\begin{enumerate}[resume]
\item Why compare medians across runs rather than means?
\item What does the permutation test assume, and what does it not assume?
\item What is the false-alarm rate at every $n$, and why is it near the test's level?
\item How many runs per side detect the planted regression 99\% of the time?
\end{enumerate}

\medskip\noindent\textbf{Part III --- Power.}
\begin{enumerate}[resume]
\item Read the detection rate of a 5\% rise at 80, 120 and 160 runs per side.
\item Interpolate the number of runs for 90\% power.
\item Compare with the normal approximation, and explain the difference.
\item How long does a gate of that size take on this laptop, at about 30 milliseconds a run?
\end{enumerate}

\medskip\noindent\textbf{Part IV --- The verdict.}
\begin{enumerate}[resume]
\item State the \emph{named result}: the number of interleaved runs the gate needs to detect a 5\% rise in the 99th percentile
with 90\% power at this machine's measured noise.
\item What would isolated cores (chapter~13) change, and why?
\item Why is a tolerance needed on top of the test?
\item What would a gate on the median latency need, and why is it not enough?
\item How would you make the gate a routine step of continuous integration without making every change wait?
\item What kinds of regression can no benchmark of the engine alone catch?
\item Why keep the pools of runs as committed data?
\item In one sentence: what makes a performance gate trustworthy?
\end{enumerate}
\end{problem}

\section{Interview questions}

\begin{interviewq}[$\star$ \iqroles{developer}]\label{iq:ll:testing-and-deploying-low-latency-systems:1}
What is a \omterm{def:ll:testing-and-deploying-low-latency-systems:tests}{property-based test}? Give an example for an order book.
\end{interviewq}

\begin{interviewq}[$\star\star$ \iqroles{developer}]\label{iq:ll:testing-and-deploying-low-latency-systems:2}
How would you fuzz a market-data parser, and what would you check besides crashes?
\end{interviewq}

\begin{interviewq}[$\star\star$ \iqroles{developer}]\label{iq:ll:testing-and-deploying-low-latency-systems:3}
Your benchmark says the new build is 3\% slower at the 99th percentile. How do you decide whether to believe it?
\end{interviewq}

\begin{interviewq}[$\star\star$ \iqroles{developer}]\label{iq:ll:testing-and-deploying-low-latency-systems:4}
What does an exchange's certification test, and what does it not?
\end{interviewq}

\begin{interviewq}[$\star\star$ \iqroles{developer, trader}]\label{iq:ll:testing-and-deploying-low-latency-systems:5}
How would you roll out a new version of a strategy engine, and how would you roll it back?
\end{interviewq}

\begin{interviewq}[$\star\star\star$ \iqroles{developer}]\label{iq:ll:testing-and-deploying-low-latency-systems:6}
Design the test and release pipeline for a low-latency trading system, from a commit to production, and say what each stage
catches.
\end{interviewq}
